\scanpage{175}
\mspage{124}
% archive sheet number: 175
Considérons dans $\unsure{\mathscr{I}|\Phi}$ les deux \deleted{ouverts}\inserted{fermés} \inserted{isom. à $\Delta^1$}

\[
\begin{aligned}
\deleted{\unsure{U}}F_1&={}_1\backslash\mathscr{I}_{\unsure{/1|\deleted{/1}}}=\{0,1\}\\
\deleted{U}F_2&={}_2\backslash\mathscr{I}_{/2}=\{0,2\},
\end{aligned}
\]
on a

\[
\deleted{U}F_1\cap\deleted{U}F_2={}_0\backslash\mathscr{I}_{\deleted{\unsure{/0}}}=\{0\},\qquad
\deleted{U}F_1\cup\deleted{U}F_2=\mathscr{I}.
\]

Soient

\[
\begin{aligned}
\deleted{\unsure{\mathscr{X}|\mathcal X}_1={}}\,{}_1\backslash\mathscr{X}_{\deleted{/1}}&=\mathscr{X}\times_{\mathscr{I}}{}_1\backslash\mathscr{I}_{\deleted{/1}}=\mathscr{X}\times_{\mathscr{I}}\deleted{U}F_1\\
\deleted{\mathscr{X}_2={}}\,{}_2\backslash\mathscr{X}_{\deleted{/2}}&=\mathscr{X}\times_{\mathscr{I}}{}_2\backslash\mathscr{I}_{\deleted{/2}}=\mathscr{X}\times_{\mathscr{I}}\deleted{U}F_2
\end{aligned}
\]
les images inverses de $\deleted{U}F_1,\deleted{U}F_2$ dans $\mathscr{X}$, on a donc

\[
{}_1\backslash\mathscr{X}_{\deleted{/1}}\cap{}_2\backslash\mathscr{X}_{\deleted{/2}}=X_0,\qquad
{}_1\backslash\mathscr{X}_{\deleted{/1}}\cup{}_2\backslash\mathscr{X}_{\deleted{/2}}=\mathscr{X},
\]

\unsure{donc|d'où} le diagramme

\diagram{\(\text{($**$)}\qquad\begin{tikzcd}[ampersand replacement=\&,column sep=small]
\& \deleted{\unsure{\mathscr{X}}}\,X_0\arrow[dl,hook,"\alpha'_1"']\arrow[dr,hook,"\alpha'_2"] \&\\
{}_1\backslash\mathscr{X}_{\deleted{/1}}\arrow[dr,hook,"i'_1"'] \& \& {}_2\backslash\mathscr{X}_{\deleted{/2}}\arrow[dl,hook,"i'_2"]\\
\& \mathscr{X} \& {,}
\end{tikzcd}\)}
déduit du diagramme similaire \deleted{\unsure{des}}

\diagram{\begin{tikzcd}[ampersand replacement=\&,column sep=small]
\& \{0\}\arrow[dl,hook,no head]\arrow[dr,hook,no head] \&\\
\deleted{U}F_1\arrow[dr,hook] \& \& \deleted{U}F_2\arrow[dl,hook]\\
\& \mathscr{I} \& {,}
\end{tikzcd}}
est (comme ce dernier) un \deleted{diagramme}\inserted{carré} d'immersion\unsure{s} \deleted{\unsure{ouvertes|fermées}}\inserted{fermées} à la fois cartésien et cocartésien. D'autre part, on a un \unsure{hom}
